\chapter{当前结论}

\section{已经实现的能力}

\begin{enumerate}
  \item 已建成 ALOHA 专用双臂数据采集软件和通用机器人数据编辑中心；
  \item 已形成多项目、多条件、多视角的真实示教数据资产；
  \item 已完成正式视觉—语言—动作模型的全参数训练、保存和现场部署；
  \item 已在现场观察到取瓶、双臂协同拧盖、瓶身与瓶盖分箱的完整循环；
  \item 已观察到长时间连续运行和一定随机条件适应表现；
  \item 已建立真实注意力记录，可观察模型使用三路视野的情况；
  \item 已开展强化学习研究，并形成数据、训练、保存和离线验证闭环。
\end{enumerate}

\section{部分实现的能力}

\begin{itemize}
  \item \textbf{泛化：}已有多种方向、无盖、翻转和多瓶数据，现场也有初步适应迹象，但没有冻结测试集与分条件成功率；
  \item \textbf{长程稳定性：}现场估计连续工作超过 1 小时，但没有起止记录、逐瓶计数和中断分类；
  \item \textbf{失败解释：}无盖指令风险有直接数据支持，倒抓仍缺逐姿态证据；
  \item \textbf{强化学习：}离线动作误差改善，但没有证明真机优于基础模型；
  \item \textbf{注意力解释：}真实记录存在，但无因果干预与结果标签。
\end{itemize}

\section{尚未实现或尚不能确认的能力}

\begin{itemize}
  \item 自动判断瓶盖是否存在、是否完全脱离以及是否投放到正确回收盒；
  \item 标准化完整任务成功率、分阶段成功率、失败频次和置信区间；
  \item 对无盖、正反方向、密度、瓶型和光照的同条件对照结果；
  \item 接触力、瓶盖转角和毫米级插管误差的测量闭环；
  \item 多随机种子训练和多个保存点的统一评测；
  \item 强化学习模型的同条件真机增益；
  \item 稳定、可自动恢复的无人值守演示。
\end{itemize}

\section{当前最好结果与主要瓶颈}

当前最好结果是：机器人在现场观察中已经能够完成完整双臂瓶子分拣循环，并表现出超过 1 小时、约 2 瓶/分钟的连续运行能力；二者均为估计而非正式测试。最主要瓶颈不是“模型完全不会做”，而是\textbf{缺少可信评测与条件判断}：没有自动成功判定，也没有逐次计数；无盖数据仍给出无条件拧盖指令；倒抓没有逐姿态诊断。

\resultbox{系统具备现场展示基础，但在没有固定测试、逐次记录和失败恢复统计之前，不宜对外宣称“稳定无人值守”或给出正式成功率。}
